% ============================================================
% CAEM — Final Research Plan (Revised April 2026)
% Author: Aksan Gony Alif, BRAC University
% Status: Post mini-run (EXP-19), cloud GPU confirmed, device unrestricted
% This document supersedes the experimental design sections of
% caem-unified-plan-v3.tex wherever they conflict.
% ============================================================

\documentclass[12pt,a4paper]{article}
\usepackage[margin=1in]{geometry}
\usepackage{amsmath,amssymb}
\usepackage{booktabs}
\usepackage{array}
\usepackage{enumitem}
\usepackage{xcolor}
\usepackage{hyperref}
\usepackage{titlesec}
\usepackage{mdframed}
\usepackage{multirow}
\usepackage{caption}

\definecolor{fixred}{RGB}{180,30,30}
\definecolor{okgreen}{RGB}{20,120,50}
\definecolor{warnblue}{RGB}{20,60,160}
\definecolor{boxbg}{RGB}{245,248,255}

\mdfdefinestyle{infobox}{
  backgroundcolor=boxbg,
  linecolor=warnblue,
  linewidth=1.2pt,
  leftmargin=0pt, rightmargin=0pt,
  innerleftmargin=10pt, innerrightmargin=10pt,
  innertopmargin=8pt, innerbottommargin=8pt,
  roundcorner=4pt
}
\mdfdefinestyle{fixbox}{
  backgroundcolor={red!5},
  linecolor=fixred,
  linewidth=1.4pt,
  leftmargin=0pt, rightmargin=0pt,
  innerleftmargin=10pt, innerrightmargin=10pt,
  innertopmargin=8pt, innerbottommargin=8pt,
  roundcorner=4pt
}

\newcommand{\caem}{CAEM}
\newcommand{\upre}{u_{\text{pre}}}
\newcommand{\uhat}{\hat{u}}
\newcommand{\ustored}{\hat{u}_{\text{stored}}}

\title{\textbf{CAEM: Final Research Plan}\\[0.4em]
\large Confidence-Aware Episodic Memory with Self-Improvement\\[0.2em]
\normalsize Revised April 2026 --- Post EXP-19, Cloud GPU Confirmed}
\author{Aksan Gony Alif\\
BRAC University\\
Supervisor: Prof.\ Md.\ Golam Rabiul Alam}
\date{April 2026}

\begin{document}
\maketitle
\tableofcontents
\newpage

% ============================================================
\section{Status Summary}
% ============================================================

\begin{mdframed}[style=infobox]
\textbf{What is complete:}
\begin{itemize}[noitemsep]
  \item Chapters 1--2: submitted (Phase 1 pre-thesis report).
  \item Chapter 3 (Requirements, Impacts, Project Management): written and compiled.
  \item Chapter 4 (Methodology): written, all 5 algorithms in \texttt{algorithmic} style,
        all figures placed, notation table, theoretical analysis including Data Purity Theorem.
  \item Implementation: all 8 pipeline stages, 96 unit tests passing, mini-run (EXP-19)
        completed on RTX~3060 (n=500, 4 cycles).
  \item Mini-run results confirmed: FEVER +10.9\%, StrategyQA +9.4\%,
        TruthfulQA net +4.0\%, HotpotQA EM$\approx$0 (boundary condition).
\end{itemize}

\textbf{What is pending:}
\begin{itemize}[noitemsep]
  \item FEVER data leakage fix (Section~\ref{sec:fever-fix}) --- \textcolor{fixred}{\textbf{blocking}}.
  \item Dataset suite revision (Section~\ref{sec:datasets}).
  \item Full-scale cloud experiment (Section~\ref{sec:experiment}).
  \item Ablation suite (Section~\ref{sec:ablations}).
  \item Chapters 5 and 6 (Section~\ref{sec:writing}).
\end{itemize}
\end{mdframed}

% ============================================================
\section{Critical Fix: FEVER Data Leakage}
\label{sec:fever-fix}
% ============================================================

\begin{mdframed}[style=fixbox]
\textbf{\textcolor{fixred}{Data integrity issue --- fix before any full-scale run.}}

\textbf{Problem:} \texttt{run\_experiment.py} loads FEVER from the \texttt{train} split
of \texttt{lucadiliello/fever} for both the self-improvement loop (SIL) query pool
\emph{and} evaluation.
Even though \texttt{split\_calibration\_sets()} carves out calibration and purity
subsets, the evaluation queries are drawn from the same \texttt{train} split.
The model is therefore evaluated on queries from the same distribution it was
trained on --- a direct data leakage.
Mini-run FEVER results (+10.9\%) are \emph{potentially inflated} and cannot be
reported as scientifically valid until the fix is applied and results re-checked.

\textbf{Root cause:} \texttt{lucadiliello/fever} has three splits:
\begin{itemize}[noitemsep]
  \item \texttt{train} --- 145,449 examples (labels available).
  \item \texttt{paper\_dev} --- 19,998 examples (labels available) --- the standard
        published benchmark split used by every FEVER baseline (DPR, FiD, Atlas).
  \item \texttt{paper\_test} --- unlabelled.
\end{itemize}
\texttt{paper\_dev} was never referenced in \texttt{run\_experiment.py}.

\textbf{Fix:}
\begin{enumerate}[noitemsep]
  \item Evaluation for FEVER must use \texttt{paper\_dev} exclusively.
        Never pass \texttt{paper\_dev} examples to the SIL query pool.
  \item \texttt{train} is allocated as follows (in order, no overlap):
    \begin{itemize}[noitemsep]
      \item Indices 0--499: calibration set (for $\uhat$ logistic regression).
      \item Indices 500--999: purity validation set (for verifier accuracy $\alpha$).
      \item Indices 1000+: SIL query pool (fed into the self-improvement loop).
    \end{itemize}
  \item Update \texttt{scripts/run\_experiment.py} and
        \texttt{caem/benchmarks/fever.py} to load \texttt{paper\_dev} for evaluation.
  \item Rerun mini-run smoke test with fix applied to confirm FEVER results are
        still positive (expected: similar magnitude, but now scientifically clean).
\end{enumerate}

\textbf{Secondary benefit:} Using \texttt{paper\_dev} as the eval set makes
CAEM's FEVER numbers directly comparable to every published baseline
that also uses \texttt{paper\_dev}, including the Liu et al.\ (ACL~2024)
RA-ISF numbers cited in PUB-01.
\end{mdframed}

% ============================================================
\section{Final Dataset Suite}
\label{sec:datasets}
% ============================================================

\subsection{Removed Datasets}

\textbf{HotpotQA --- removed entirely.}
EXP-19 confirmed Flan-T5-Large achieves EM~$\approx$~0.006 on HotpotQA across all
four cycles.
The Data Purity Theorem's precondition $p > (1-\alpha)$ is not satisfied when
base accuracy is near zero.
This result is reported as a boundary condition validation in Chapter~5, using
mini-run data (no further HotpotQA compute needed).
Do not include HotpotQA in the full-scale experiment.

\subsection{Training Datasets (Self-Improvement Loop)}

These datasets feed the SIL query pool across all cycles.
Selection criteria: (a)~sufficiently large training split for 10 cycles at 5,000
queries/cycle without repetition; (b)~Flan-T5-Large has non-trivial base accuracy
($p \gg 0$), satisfying the purity theorem precondition;
(c)~verifiable answers compatible with NLI + SC + SE signals.

\begin{table}[htbp]
\centering
\caption{Self-improvement loop training datasets (final).}
\label{tab:train-datasets}
\renewcommand{\arraystretch}{1.3}
\begin{tabular}{@{}p{2.8cm}p{1.8cm}p{2.0cm}p{6.0cm}@{}}
\toprule
\textbf{Dataset} & \textbf{Train size} & \textbf{Cycles covered} & \textbf{Rationale} \\
\midrule
FEVER & 145,449 & 29$\times$ & Fact verification; NLI maps directly onto task;
  primary verifier signal. Fix leakage before use (Section~\ref{sec:fever-fix}). \\
TriviaQA & 78,785 & 15$\times$ & Factoid; diverse domains; Flan-T5-Large has
  reasonable base capability. Unique in testing parametric knowledge breadth
  rather than reasoning chains. \\
Natural Questions & 79,168 & 15$\times$ & Real user queries; open-domain;
  large; diverse. Complements FEVER (verification) and TriviaQA (trivia)
  with grounded natural-language questions. \\
\bottomrule
\end{tabular}
\end{table}

Each training dataset provides 5,000 queries per cycle.
At 10 cycles, each dataset contributes 50,000 total queries --- well within its
training split without repetition.

\textbf{Why not StrategyQA as training?}
StrategyQA has only 2,290 training samples.
At 5,000 queries/cycle, the pool is exhausted before Cycle~1 completes.
Repeated passes risk data memorisation rather than self-improvement.
StrategyQA is re-assigned to the transfer evaluation suite
(Section~\ref{sec:eval-datasets}).

\subsection{Evaluation Datasets (Transfer / Held-Out)}
\label{sec:eval-datasets}

These datasets are \emph{never} seen during training or the SIL loop.
They are evaluated post-cycle to measure out-of-distribution generalisation.

\begin{table}[htbp]
\centering
\caption{Transfer and held-out evaluation datasets (final).}
\label{tab:eval-datasets}
\renewcommand{\arraystretch}{1.3}
\begin{tabular}{@{}p{2.8cm}p{1.6cm}p{6.8cm}@{}}
\toprule
\textbf{Dataset} & \textbf{Eval size} & \textbf{Role and rationale} \\
\midrule
TruthfulQA & 817 & Adversarial misconceptions; tests whether hallucination
  reduction on factoid/verification tasks transfers to misconception resistance.
  Evaluation only --- no training. Primary transfer claim. \\
StrategyQA & 490 (test) & Implicit multi-step reasoning; tests transfer to
  reasoning tasks not in the training distribution. \\
ARC-Challenge & 1,172 & Science reasoning; hardest tier of ARC benchmark;
  tests transfer to domain-specific factual inference beyond Wikipedia. \\
HotpotQA (boundary) & 500 & Mini-run data only. Reports theorem boundary
  condition ($p \approx 0 \Rightarrow$ no improvement). No new compute. \\
\bottomrule
\end{tabular}
\end{table}

\textbf{Transfer narrative:}
If \caem{} trained on FEVER + TriviaQA + Natural Questions improves on
TruthfulQA, StrategyQA, and ARC-Challenge --- benchmarks it never trained on ---
the system demonstrates generalised hallucination reduction, not just
in-distribution optimisation.
This is a stronger claim than improving on the training benchmarks alone.

% ============================================================
\section{System Configuration (Final)}
\label{sec:config}
% ============================================================

\begin{table}[htbp]
\centering
\caption{Final experimental configuration. All values are design decisions
(\texttt{[DES]}) unless marked \texttt{[LIT]} (literature-fixed) or
\texttt{[CAL]} (empirically calibrated).}
\label{tab:config}
\renewcommand{\arraystretch}{1.35}
\begin{tabular}{@{}p{5.0cm}p{3.5cm}p{5.0cm}@{}}
\toprule
\textbf{Parameter} & \textbf{Value} & \textbf{Rationale} \\
\midrule
Base model & Flan-T5-Large (780M) \texttt{[LIT]} &
  Unchanged. Instruction-tuned; multi-task capable; data purity theorem
  boundary condition validated on this scale. \\
\# self-improvement cycles & 10 \texttt{[DES]} &
  Theory predicts practical equilibrium at 7--10 cycles.
  10 cycles captures full convergence curves for all three training benchmarks. \\
Queries per benchmark per cycle & 5,000 \texttt{[DES]} &
  Sufficient for statistically meaningful per-cycle EM; no dataset
  exhaustion across 10 cycles given training split sizes. \\
Max episodic memory & 1,000,000 episodes \texttt{[DES]} &
  Removes capacity as a confounding variable. Pruning trigger at 95\%
  (950,000) retained as safety ceiling. FAISS IVF-PQ at 1M scale:
  $\sim$1.5\,GB for compressed vectors. \\
FAISS index type (memory) & IVF-PQ, \texttt{nlist}=4096,
  $M$=64, \texttt{nbits}=8 \texttt{[DES]} &
  Handles 1M-scale retrieval efficiently within A100 VRAM. \\
RAG passage corpus & Full DPR Wikipedia \texttt{[LIT]} &
  21,015,324 passages (100-word chunks, Karpukhin et al.\ 2020).
  Standard split used by all published RAG baselines (DPR, FiD, Atlas) ---
  ensures comparability. Needs $\geq$128\,GB system RAM for CPU FAISS index. \\
RAG FAISS index type & IVF-PQ, \texttt{nlist}=65536 \texttt{[DES]} &
  Scales to 21M passages; $\sim$40\,GB compressed. \\
Embedding model & \texttt{all-mpnet-base-v2} \texttt{[LIT]} &
  384-dim output. Unchanged. \\
Verification NLI model & RoBERTa-Large-MNLI \texttt{[LIT]} & Unchanged. \\
L2 regularisation $\lambda$ & 0.01 \texttt{[DES]} &
  Adapted from Kirkpatrick et al.\ (2017); validated in EWC ablation
  (Section~\ref{sec:ablations}). \\
Fine-tuning mix & 90\% verified / 10\% general \texttt{[DES]} & Unchanged. \\
Forgetting abort threshold & Retention $< 0.93$ per cycle \texttt{[DES]} &
  Unchanged. Checked against MMLU + HellaSwag + BoolQ + GSM8K. \\
$\uhat$ accept threshold & 0.60 \texttt{[DES]} & Unchanged. \\
Tier~1 routing threshold & 0.90 \texttt{[DES]} & Unchanged. \\
OR-condition override & $\upre < 0.60$ \texttt{[DES]} & Unchanged. \\
MC Dropout passes ($\uhat$) & $K=5$ \texttt{[LIT]} & Unchanged. \\
Self-consistency chains & $M=3$ \texttt{[LIT]} & Unchanged. \\
Semantic entropy samples & $K=10$ \texttt{[LIT]} & Unchanged. \\
Training device & A100 80\,GB SXM4 (cloud) & Full fp32 for both models;
  batch size 16; $\sim$6--8\,h inference/cycle + 20--30\,min FT/cycle. \\
\bottomrule
\end{tabular}
\end{table}

\subsection{Memory at 1M Scale --- Practical Notes}

\textbf{FAISS IVF-PQ at 1M episodes:}
Compressed vectors at $M$=64 bytes/vector: $1{,}000{,}000 \times 64 = 64$\,MB.
Centroid table: $4096 \times 384 \times 4 = 6$\,MB.
Total in-memory footprint: $\approx$70\,MB --- trivially within A100 VRAM
alongside both models ($\sim$4\,GB fp32 total).

\textbf{Retroactive re-verification cost at 1M:}
Re-verifying all stored episodes at the end of each cycle scales linearly.
Apply the delta-threshold optimisation: only re-verify episodes where
$\ustored \in [0.45, 0.70]$ (the uncertainty band).
Episodes with $\ustored > 0.70$ are stable; episodes with $\ustored < 0.45$
are pruning candidates regardless.
This reduces the re-verification pass to $\approx$30--40\% of stored episodes
per cycle.

\textbf{Full Wikipedia DPR index:}
The standard DPR Wikipedia index (21M passages) is available as
\texttt{wiki\_dpr} on HuggingFace or via the original DPR repository.
Build once; the index is static across all cycles.
Requires $\geq$128\,GB system RAM for CPU-resident FAISS search during inference
(passage retrieval does not require GPU).
Cloud instances with 192\,GB RAM (e.g.\ \texttt{a2-highgpu-1g} on GCP or
\texttt{p3.8xlarge} with CPU offload on AWS) are adequate.

% ============================================================
\section{Experimental Design}
\label{sec:experiment}
% ============================================================

\subsection{Data Splits (Non-Overlapping --- Critical)}

\begin{table}[htbp]
\centering
\caption{Data split allocation. Splits must be non-overlapping across all uses.}
\label{tab:splits}
\renewcommand{\arraystretch}{1.3}
\begin{tabular}{@{}p{3.2cm}p{2.4cm}p{2.4cm}p{5.2cm}@{}}
\toprule
\textbf{Purpose} & \textbf{Size} & \textbf{Source split} & \textbf{Notes} \\
\midrule
\multicolumn{4}{l}{\textit{FEVER}} \\
Calibration & 500 & \texttt{train}[0:500] &
  Fits $\uhat$ logistic regression weights. \\
Purity validation & 500 & \texttt{train}[500:1000] &
  Measures verifier accuracy $\alpha$. \\
SIL query pool & train[1000+] & \texttt{train}[1000:] &
  Fed into self-improvement loop. \\
Evaluation & 19,998 & \texttt{paper\_dev} &
  Never seen during training. Standard benchmark split. \\
\midrule
\multicolumn{4}{l}{\textit{TriviaQA}} \\
Calibration & 500 & \texttt{train}[0:500] & \\
Purity validation & 500 & \texttt{train}[500:1000] & \\
SIL query pool & train[1000+] & \texttt{train}[1000:] & \\
Evaluation & \texttt{validation} split & \texttt{validation} &
  Official held-out validation set. \\
\midrule
\multicolumn{4}{l}{\textit{Natural Questions}} \\
Calibration & 500 & \texttt{train}[0:500] & \\
Purity validation & 500 & \texttt{train}[500:1000] & \\
SIL query pool & train[1000+] & \texttt{train}[1000:] & \\
Evaluation & \texttt{validation} split & \texttt{validation} &
  Official held-out validation set. \\
\midrule
\multicolumn{4}{l}{\textit{Transfer evaluation (no training split used)} } \\
TruthfulQA & 817 & \texttt{validation} & Entire dataset for eval. \\
StrategyQA & 490 & \texttt{test} & Official test split. \\
ARC-Challenge & 1,172 & \texttt{test} & Official test split. \\
\bottomrule
\end{tabular}
\end{table}

\subsection{Self-Improvement Loop Execution}

\begin{enumerate}
  \item \textbf{Cold-start seeding:} Seed episodic memory with $\sim$200 verified
        episodes per training benchmark (600 total) from the SIL query pool
        before Cycle~0.
        This prevents empty-memory degradation in Cycle~1.
  \item \textbf{Per-cycle loop} (repeat for cycles 0--9):
    \begin{enumerate}[label=\alph*.]
      \item Sample 5,000 queries per training benchmark from the SIL pool
            (without replacement across cycles).
      \item Run each query through the full 8-stage pipeline
            (encode $\to$ retrieve $\to$ confidence $\to$ route $\to$ generate
            $\to$ verify $\to$ store).
      \item Collect verified episodes into episodic memory.
      \item At end of cycle: retroactive re-verification of uncertainty-band
            episodes ($\ustored \in [0.45, 0.70]$).
      \item Fine-tune on verified episodes (90/10 mix, L2 regularisation,
            $\lambda = 0.01$).
      \item Check forgetting abort: if MMLU/HellaSwag/BoolQ/GSM8K retention
            $< 0.93$, revert to previous cycle weights.
      \item Evaluate on all benchmark eval splits (training + transfer).
            Save per-question outputs for ablation analysis.
    \end{enumerate}
  \item \textbf{Calibration:} After Cycle~0 completes, fit $\uhat$ logistic
        regression weights on the 500-sample calibration sets.
        Update \texttt{calibration/calibrated\_config.json}.
        Re-run Cycles~1--9 with calibrated weights.
\end{enumerate}

\subsection{Baselines}

All baselines are evaluated once (no self-improvement loop) and reported in
Table~5.1 (main results).

\begin{table}[htbp]
\centering
\caption{Baseline methods for Chapter~5 comparison.}
\label{tab:baselines}
\renewcommand{\arraystretch}{1.25}
\begin{tabular}{@{}p{3.8cm}p{9.0cm}@{}}
\toprule
\textbf{Baseline} & \textbf{What it lacks vs.\ \caem{}} \\
\midrule
Zero-shot Flan-T5-Large & No memory, no verification, no self-improvement. \\
Chain-of-Thought (3-shot) & No memory, no verification, no self-improvement. \\
RAG (full Wikipedia) & No memory accumulation, no confidence routing,
  no self-improvement. \\
Self-Consistency (M=3) & No memory, no routing, no self-improvement.
  Verification only. \\
Vanilla Fine-Tuning & No memory, no confidence gating, no retroactive cleaning.
  Trains on all generated data regardless of quality. \\
Memory-Only (no FT) & No fine-tuning. Tests whether memory accumulation alone
  (without model improvement) explains gains. \\
GPT-3.5 vanilla & Published numbers from Liu et al.\ ACL Findings 2024
  (RA-ISF, arXiv:2403.06840). No compute required. \\
GPT-3.5 + RAG & Published numbers from same source. \\
Self-RAG 13B & Published numbers from same source.
  Key comparison: CAEM at 780M vs.\ Self-RAG at 13B. \\
\bottomrule
\end{tabular}
\end{table}

% ============================================================
\section{Ablation Studies}
\label{sec:ablations}
% ============================================================

\subsection{Internal Ablations (run\_ablation.py --- already implemented)}

Nine configurations grouped by the mechanism they probe.
Run after the full-scale experiment using saved per-question outputs.

\begin{table}[htbp]
\centering
\caption{Internal ablation configurations (Chapter~5, Table~5.5).}
\label{tab:ablations}
\renewcommand{\arraystretch}{1.3}
\begin{tabular}{@{}p{1.2cm}p{3.5cm}p{8.0cm}@{}}
\toprule
\textbf{ID} & \textbf{Config} & \textbf{Committee question answered} \\
\midrule
\multicolumn{3}{l}{\textit{Memory quality group}} \\
A1 & No retroactive re-verification &
  Is memory quality improvement bidirectional, or does it only flow forward? \\
A-NLI & Single-layer NLI only &
  Do SC and SE add value beyond NLI alone? \\
A-verify & No verification (store all) &
  Does the verification gate prevent memory contamination? \\
\midrule
\multicolumn{3}{l}{\textit{Routing group}} \\
A2 & All-Tier-3 (no routing) &
  Does the routing architecture provide a computational dividend? \\
A3 & No OR-condition override &
  Does the safety override prevent routing failures? \\
\midrule
\multicolumn{3}{l}{\textit{Training group}} \\
A4 & Vanilla fine-tuning (no L2) &
  Does L2 regularisation prevent catastrophic forgetting? \\
A5 & No fine-tuning (memory only) &
  Is fine-tuning necessary, or does memory accumulation alone explain gains? \\
\midrule
\multicolumn{3}{l}{\textit{Confidence group}} \\
A6 & Remove semantic entropy from $\uhat$ &
  Does SE uniquely contribute to hallucination detection? \\
A7 & $(q, a)$-only training (no reasoning chain $c$) &
  Does reasoning-chain supervision improve over answer-only supervision? \\
\bottomrule
\end{tabular}
\end{table}

\subsection{Extended Ablations (cloud GPU --- new experiments)}

These were previously compute-blocked on RTX~3060.
All are now feasible on A100~80\,GB.
Merge PUB-04 and PUB-05 into a single 4-variant fine-tuning run to avoid
running two separate training sweeps.

\begin{table}[htbp]
\centering
\caption{Extended ablations (new experiments on cloud GPU).}
\label{tab:pub-ablations}
\renewcommand{\arraystretch}{1.3}
\begin{tabular}{@{}p{1.4cm}p{3.0cm}p{2.0cm}p{6.8cm}@{}}
\toprule
\textbf{ID} & \textbf{Config} & \textbf{Est.\ cost} & \textbf{Purpose} \\
\midrule
PUB-04+05 & Four-variant FT comparison:
  (i) full FT + L2,
  (ii) full FT + EWC,
  (iii) LoRA ($r$=8, $\alpha$=16) + L2,
  (iv) LoRA alone &
  $\sim$8\,h & Validates L2 as correct EWC approximation (PUB-05) and
  pre-empts ``why not PEFT?'' question (PUB-04). If LoRA + L2 matches
  full FT + L2 on EM with better MMLU retention, recommend as future
  deployment path. \\
PUB-06 & Flan-T5-XL (3B) on FEVER only &
  $\sim$6\,h & Tests whether \caem{} mechanism holds at larger scale.
  Also partially redeems the HotpotQA capability-ceiling finding:
  if 3B clears the $p > (1-\alpha)$ precondition on HotpotQA, confirm
  as a note. \\
PUB-03 & Human evaluation on TruthfulQA (N=100) &
  No GPU & Two annotators evaluate Cycle~0 vs.\ Cycle~10 output pairs
  blind. Report Cohen's $\kappa$. Validates ROUGE-L proxy used in
  automated evaluation. \\
\bottomrule
\end{tabular}
\end{table}

\subsection{What to Drop}

\textbf{Full-scale HotpotQA (780M) run --- dropped.}
EXP-19 already confirmed EM~$\approx$~0.006 across 4 cycles at $n$=500.
Scaling to $n$=5,000 on a benchmark the model cannot answer adds no new
information.
Report EXP-19 mini-run data as the boundary condition finding in Chapter~5.
This saves $\sim$8--10 GPU hours.

% ============================================================
\section{Labels and Scientific Validity}
\label{sec:labels}
% ============================================================

This section records the final position on ground-truth dependency ---
the answer to the supervisor's question.

\begin{table}[htbp]
\centering
\caption{Ground-truth label dependency by component.}
\label{tab:labels}
\renewcommand{\arraystretch}{1.3}
\begin{tabular}{@{}p{4.5cm}p{2.2cm}p{6.0cm}@{}}
\toprule
\textbf{Component} & \textbf{Labels needed?} & \textbf{What it uses instead} \\
\midrule
NLI verification (Layer~1) & No &
  Retrieved Wikipedia passage as premise; generated answer as hypothesis.
  Self-supervised. \\
Self-consistency (Layer~2) & No &
  $M$=3 model outputs compared against each other. \\
Semantic entropy (Layer~3) & No &
  $K$=10 samples; NLI-based equivalence clustering. \\
Self-improvement fine-tuning & No &
  Trains on model's own verified (passage, reasoning, answer) triples. \\
$\uhat$ weight calibration & $\sim$500 samples &
  Logistic regression on (signal values, correct/incorrect) pairs.
  One-time bootstrap; no ongoing label dependency. \\
Cold-start seeding & Optional &
  Can use existing QA datasets or self-generate + verify seed episodes. \\
Evaluation & Yes (benchmark) &
  Ground-truth labels used for EM/F1 measurement only ---
  not part of the training loop. \\
\bottomrule
\end{tabular}
\end{table}

\textbf{Supervisor response (canonical):}
``\caem{} uses approximately 500 labels per training benchmark for a one-time
confidence calibration step.
These labels calibrate \emph{how much to trust the model's confidence signal},
not what the correct answers are.
After calibration, the self-improvement loop runs indefinitely on self-generated,
self-verified data with no additional labels.
A direct fine-tune on 500 samples would give a marginally better static model
with no confidence signal, no routing, and no continued improvement.
\caem{}'s 500 labels buy a self-improving loop that runs across 10 cycles and
transfers to held-out benchmarks the model was never trained on.''

% ============================================================
\section{Writing Plan}
\label{sec:writing}
% ============================================================

\begin{table}[htbp]
\centering
\caption{Chapter writing status and dependencies.}
\label{tab:writing}
\renewcommand{\arraystretch}{1.3}
\begin{tabular}{@{}p{1.8cm}p{3.5cm}p{2.5cm}p{4.5cm}@{}}
\toprule
\textbf{Chapter} & \textbf{Content} & \textbf{Status} & \textbf{Blocking dependency} \\
\midrule
Ch~1 & Introduction & \textcolor{okgreen}{\textbf{DONE}} & None \\
Ch~2 & Literature Review & \textcolor{okgreen}{\textbf{DONE}} & None \\
Ch~3 & Requirements, Impact, Project Management &
  \textcolor{okgreen}{\textbf{DONE}} & None \\
Ch~4 & Methodology (full) & \textcolor{okgreen}{\textbf{DONE}} &
  Add formal purity theorem proof (PUB-02, $\sim$2h, no data needed) \\
Ch~5 & Experimental Results & \textcolor{fixred}{\textbf{NOT WRITTEN}} &
  Needs full-scale cloud experiment results \\
Ch~6 & Conclusion & \textcolor{fixred}{\textbf{NOT WRITTEN}} &
  Needs Ch~5 complete \\
\bottomrule
\end{tabular}
\end{table}

\subsection{Chapter 5 Structure}

\begin{enumerate}
  \item \textbf{\S5.1 Experimental Setup} --- hardware (A100 cloud), datasets
        (Table~\ref{tab:train-datasets} + Table~\ref{tab:eval-datasets}),
        baselines (Table~\ref{tab:baselines}), evaluation metrics (EM/F1,
        HallRed\%, MMLU retention).
  \item \textbf{\S5.2 Main Results} --- per-cycle accuracy table across all
        training and transfer benchmarks; primary narrative thread:
        ``Does \caem{} outperform baselines lacking a closed loop?''
  \item \textbf{\S5.3 Mechanism Evidence} --- three evidence tables:
        (a)~Tier~1\% growth per cycle (memory accumulation);
        (b)~$\bar{u}_{\text{stored}}$ rising + stored\% falling per cycle
        (memory quality); (c)~retroactive re-verification delta per cycle
        (backward purity improvement).
  \item \textbf{\S5.4 Transfer Evaluation} --- TruthfulQA, StrategyQA,
        ARC-Challenge results post-Cycle~10.
        Headline: generalised hallucination reduction without task-specific
        adaptation.
  \item \textbf{\S5.5 Theoretical Validation} --- purity theorem empirical
        check; convergence curve vs.\ theoretical prediction;
        HotpotQA boundary condition (mini-run data).
  \item \textbf{\S5.6 Ablation Study} --- internal ablations
        (Table~\ref{tab:ablations}) + extended ablations
        (Table~\ref{tab:pub-ablations}).
        Each ablation answers one specific committee question.
  \item \textbf{\S5.7 Efficiency and Stability} --- per-tier latency,
        Tier~1 fraction growth (latency dividend), MMLU retention across
        cycles, forgetting abort activations (if any), total GPU hours.
  \item \textbf{\S5.8 Error Analysis} --- HotpotQA capability ceiling
        (theory validation); TruthfulQA ROUGE-L instability; human eval
        agreement (PUB-03).
\end{enumerate}

\subsection{Chapter 6 Structure}

\begin{enumerate}
  \item \textbf{\S6.1 Summary of Findings} --- one paragraph per confirmed
        mechanism + one for the HotpotQA boundary condition.
  \item \textbf{\S6.2 Limitations} --- model scale dependency (780M ceiling);
        calibration label requirement ($\sim$500 samples); inference scope
        (NLI requires retrieved context in open-domain mode);
        ROUGE-L unreliability for TruthfulQA; single-GPU reproducibility
        claim updated to A100 cloud.
  \item \textbf{\S6.3 Future Work} --- (i)~model capacity scaling
        (Flan-T5-XL/XXL); (ii)~knowledge distillation from verified memory
        into smaller student models; (iii)~reference-free NLI using retrieved
        passages instead of benchmark labels; (iv)~LoRA fine-tuning path
        for deployment efficiency (from PUB-04+05 if it matches full FT).
\end{enumerate}

% ============================================================
\section{Execution Sequence}
\label{sec:execution}
% ============================================================

\begin{mdframed}[style=infobox]
\textbf{Parallel track A --- RTX 3060, start immediately (no cloud needed):}
\begin{enumerate}[noitemsep]
  \item Apply FEVER leakage fix to \texttt{run\_experiment.py} and
        \texttt{caem/benchmarks/fever.py}.
  \item Re-run mini smoke test ($n$=50, 2 cycles) to confirm FEVER results
        survive the fix.
  \item Run \texttt{scripts/run\_ablation.py} $\to$ MMLU retention numbers
        (needed for Ch~5 \S5.7).
  \item Run \texttt{scripts/run\_purity\_validation.py} $\to$ Theory~1/2/3
        validation tables (needed for Ch~5 \S5.5).
  \item Write formal purity theorem proof block in Chapter~4 (PUB-02,
        $\sim$2\,h, zero compute).
  \item Update \texttt{config.py}:
        \texttt{max\_memory\_size = 1\_000\_000},
        \texttt{n\_cycles = 10},
        \texttt{faiss\_nlist = 4096}.
\end{enumerate}

\textbf{Parallel track B --- Cloud A100 80\,GB setup:}
\begin{enumerate}[noitemsep]
  \item Add TriviaQA and Natural Questions benchmark adapters
        (\texttt{caem/benchmarks/triviaqa.py},
         \texttt{caem/benchmarks/natural\_questions.py}).
  \item Add StrategyQA and ARC-Challenge as transfer-eval-only adapters.
  \item Download and build full DPR Wikipedia FAISS index
        (21M passages; build once, reuse across all cycles).
  \item Run \texttt{scripts/seed\_cold\_start.py} for FEVER + TriviaQA + NQ
        ($\sim$600 total seed episodes).
\end{enumerate}

\textbf{Main cloud experiment (after tracks A + B complete):}
\begin{enumerate}[noitemsep]
  \item Full-scale run: FEVER + TriviaQA + NQ, 10 cycles, 5,000/benchmark/cycle,
        1M memory, full DPR Wikipedia.
  \item Post-cycle transfer eval: TruthfulQA, StrategyQA, ARC-Challenge
        (zero-shot, no training).
  \item PUB-04+05: 4-variant fine-tuning ablation (full FT + L2 / + EWC;
        LoRA + L2 / alone).
  \item PUB-06: Flan-T5-XL (3B) on FEVER, 3 cycles only.
  \item PUB-03: Human eval on TruthfulQA 100-sample subset (no GPU).
\end{enumerate}

\textbf{Writing (after main experiment):}
\begin{enumerate}[noitemsep]
  \item Chapter~5 (all sections, full data available).
  \item Chapter~6 (findings + limitations + future work).
  \item Appendices (extended ablation tables, XL model results,
        full per-question analysis).
\end{enumerate}
\end{mdframed}

\subsection{Critical Path}

\[
\underbrace{\text{FEVER fix}}_{\text{now}} \;\to\;
\underbrace{\text{cloud setup}}_{\text{parallel}} \;\to\;
\underbrace{\text{full experiment}}_{\text{blocking for Ch~5}} \;\to\;
\underbrace{\text{Ch~5}}_{\text{blocking for Ch~6}} \;\to\;
\underbrace{\text{Ch~6}}_{\text{submission}}
\]

The FEVER fix is the single earliest blocking item.
Do not start the full-scale cloud run until the fix is applied and
smoke-tested.

\end{document}
